% =============================================================================
% Chapter 1 — Introduction (REVISED)
% Tara Torbati — Master's Thesis
%
% Revisions in this version:
%   - §1.1 rewritten to restore the argumentative arc from water_scarcity:
%       supply paradox → rebound effect → governance failure → technical opening
%   - §1.4 (Scope and Limitations) compressed from 5-bullet list to 1 paragraph
%   - §1.2 and §1.3 tightened for consistency with the new §1.1 framing
%   - Sentence-length pass throughout: shorter sentences, fewer em-dash chains
%   - Rhetorical intensifiers stripped where they shaded into advocacy
% =============================================================================

\chapter{Introduction}
\label{chap:intro}

\section{Motivation and Context}
\label{sec:intro-motivation}

Iran's water-management system displays a counterintuitive pattern that
has become known as the supply paradox. Between 1990 and 2015, national
dam capacity tripled, while groundwater withdrawal simultaneously rose
from approximately 60 to 90~billion cubic meters per
year~\cite{Mesgaran2018,Behling2022}. Each new reservoir was, in
isolation, a defensible response to recurrent drought. In aggregate,
they failed to reduce extraction pressure on aquifers. The reason is
that supply-side infrastructure increases availability without
disciplining demand. As long as agricultural producers face no
volumetric cost for the water they use, additional supply is absorbed
by additional consumption rather than by storage or
conservation~\cite{Mesgaran2018,Nouri2023}.

Two structural features make this dynamic difficult to reverse through
demand-side policy. The first is the decoupling of price from use.
Iranian agricultural water tariffs are set at levels two to three
orders of magnitude below domestic and industrial tariffs, and energy
for groundwater pumping is heavily subsidized. The marginal cost of
water at the farmgate is therefore close to zero, and the price signal
that would otherwise discipline consumption is absent. The second is
the proliferation of informal wells. A significant share of
agricultural extraction occurs through wells that are unregistered,
unmetered, or only weakly monitored. Without reliable measurement,
volumetric quotas and extraction limits cannot be
enforced~\cite{Mesgaran2018,PourgholamAmiji2024}.

A natural intuition is that improving irrigation efficiency at the
field level would reduce aggregate demand. Empirically, this has not
occurred. The phenomenon has been studied under the name of the
rebound effect: when farmers adopt sprinkler or drip systems that
deliver more water to the crop per unit applied, they typically
respond by expanding cultivated area, intensifying cropping, or
shifting to more water-demanding crops. The result is that total
basin-level water consumption rises rather than
falls~\cite{PerezBlanco2020,Hesabi2023}. The mechanism is not
behavioral irrationality. It is a predictable response to volumetric
scarcity that is not bound at the system level: efficiency gains
relax the local constraint on yield, and the released capacity flows
into expansion. Conventional policy instruments intended to constrain
expansion --- volumetric quotas, withdrawal limits, regulatory
restrictions --- run into the governance limitations described
above~\cite{PerezBlanco2020,Hajirad2025}.

The combination of these dynamics defines the operational regime in
which any technical intervention must function. Supply expansion does
not discipline demand. Price signals do not reach the user.
Efficiency improvements rebound into expansion. Policy reform alone
has been observed to be insufficient. Technology adoption alone has
been observed to be insufficient. What remains viable is a
combination: setting volumetric allocations externally, through
aquifer-level monitoring or regional allocation policy, and enforcing
those allocations at the field level through a controller that
respects them as hard constraints. The technical question this thesis
addresses is whether such a controller can deliver agronomically
acceptable yields under hard volumetric scarcity, and how the
available water should be allocated across space and time on a
topographically heterogeneous field.

This question is examined in the specific context of Hashemi rice
cultivation in Gilan Province on Iran's Caspian coast. Gilan is the
country's principal rice-producing region, and rice is among the most
water-intensive crops grown
nationally~\cite{jalali2021,sadidi2021,surfiran2025}. The province
exemplifies the supply-paradox dynamics described
above~\cite{Mesgaran2018,PourgholamAmiji2024}. A controller that
performs well in this setting --- under deficit allocations, on
realistic topography, on a high-fidelity crop-soil model ---
demonstrates the feasibility of the more general proposition that
constrained technical control can substitute for the missing economic
and policy constraints on water use.

\section{Problem Statement}
\label{sec:intro-problem}

This thesis develops and compares constrained optimal-control and
reinforcement-learning approaches to irrigation on a 130-agent
topographically heterogeneous Hashemi rice field in Gilan Province,
Iran. The field is represented as a discrete grid of cells with a
non-trivial elevation gradient (111~m of vertical relief over the
field extent). Each cell maintains its own coupled soil-moisture,
phenological, and biomass states. Cells exchange water laterally
through topographically driven surface runoff. Each cell admits an
independent irrigation control input bounded by the physical actuator
capacity. The joint control vector is subject to a global
season-cumulative water budget representing the realistic deficit
allocations issued under regional water-management
schemes~\cite{Mesgaran2018,Nouri2023}.

Within this setting, the thesis investigates four questions:

\begin{enumerate}
\item \textbf{Cost-function design.} How should the multi-objective
cost function of a constrained MPC controller be structured and
weighted to balance terminal biomass, water consumption, drought
risk, and waterlogging avoidance? How should those weights be
calibrated against real economic price signals?

\item \textbf{Pathology characterization.} What dynamic pathologies
arise in the controller-plant interaction at long prediction
horizons, in saturated-soil regimes, and under cheap water? Can
these be resolved through cost-function design rather than through
ad-hoc heuristic correction?

\item \textbf{MPC versus RL comparison.} Under a fair, exact-equivalent
reward formulation, does a Soft Actor-Critic agent trained on the
ABM reproduce the policy quality of constrained MPC? At what
trade-off between training cost and deployment latency?

\item \textbf{Robustness to scenario and budget.} How do both
controllers respond to climatic scenario variation (dry, moderate,
wet) and to deficit budget allocations (100\%, 85\%, 70\% of the
climatologically derived full requirement)?
\end{enumerate}

The answers to these questions are intended to inform two distinct
audiences. For control-systems researchers, the thesis offers a
concrete case study in constrained distributed-parameter control
under economic anchoring. For water-management practitioners
considering technical deployment options under volumetric
deficit-irrigation policy, the thesis quantifies the agronomic
performance achievable under different scarcity regimes.

\section{Contributions}
\label{sec:intro-contributions}

The principal contributions of this thesis are as follows.

\paragraph{A topographical agent-based virtual plant.}
A 130-agent ABM of a Gilan rice field is developed, parameterized from
the FAO Irrigation and Drainage Paper No.~56~\cite{allen1998},
regional agronomic literature~\cite{jalali2021,sadidi2021}, and the
soil-hydraulic database of Rawls et al.~\cite{rawls1982}. The model
extends the topographical irrigation framework of Lopez-Jimenez
et~al.~\cite{Lopezjimenez2024,LopezJimenez2022} with three
modifications required for closed-loop optimal control: explicit
two-layer hydrological decoupling between surface ponding and
root-zone moisture, a padded digital-elevation-model boundary
condition that resolves a degenerate sink behavior in the original
formulation, and a coupling of drought stress directly to biomass
accumulation. The plant constitutes the shared simulation environment
on which both predictive and learning-based controllers are evaluated.

\paragraph{A constrained MPC controller with economically anchored
cost weights.} A nonlinear MPC formulation is developed for the ABM,
solved via the IPOPT interior-point method~\cite{wachter2006ipopt}
within the CasADi symbolic optimization
framework~\cite{andersson2019casadi}. The cost function is calibrated
to the four-tier Iranian water-pricing schedule. This calibration
allows the controller's behavior to be interpreted as a response to
real economic price signals rather than to abstract regularization
constants.

\paragraph{A systematic weight-sensitivity analysis identifying and
resolving three controller pathologies.} A thirty-three-configuration
sensitivity sweep is conducted across one-at-a-time, joint factorial,
cross-scenario validation, and ceiling phases. The sweep characterizes
three pathologies of the original cost-function formulation: silent
over-irrigation under cheap water, field-capacity overshoot at long
prediction horizons, and chronic waterlogging in wet-year scenarios.
The analysis further demonstrates that two of the originally proposed
penalty terms (surface ponding and field-capacity overshoot) are
dynamically redundant once one of them is sufficiently strong,
motivating a parsimonious five-term cost function. The recommended
operating point is identified and validated across all combinations
of climate scenario, budget allocation, and prediction horizon.

\paragraph{A direct MPC--SAC comparison under exact reward equivalence.}
A Soft Actor-Critic
agent~\cite{haarnoja2018sac,haarnoja2019sac} is trained against a
reward function constructed as the exact mathematical negation of
the MPC path cost. This design eliminates the common confound in
AI-versus-control benchmarking, in which differences in the objective
function obscure differences in the optimization mechanism. The SAC
architecture employs Centralized Training with Decentralized
Execution~\cite{amato2024ctde}, with parameter sharing across the 130
spatial agents.

\paragraph{Comparative results across scenario, budget, and horizon.}
All four controllers (no-irrigation baseline, fixed-schedule
heuristic, constrained MPC, and SAC) are evaluated on a 9-cell
factorial grid spanning three climate scenarios (dry 2022, moderate
2020, wet 2024) and three budget allocations (100\%, 85\%, 70\% of
the climatological full demand). The results identify the operational
regimes in which each architecture excels and quantify the
yield-versus-water-versus-latency trade-off as a function of
deployment context.

\section{Scope and Limitations}
\label{sec:intro-scope}

The thesis operates within a deliberately bounded scope. The
contribution is methodological: a comparative study of constrained
MPC and SAC on a physics-grounded, parameter-validated test bed in
a representative water-scarce setting. Empirical validation against
direct field observations in Gilan and the practical realization of
cell-level variable-rate irrigation through commercial hardware are
not undertaken; both are identified as natural directions for
future work in Chapter~\ref{chap:conclusion}.

\section{Thesis Structure}
\label{sec:intro-structure}

The remainder of this thesis is organized as follows.
Chapter~\ref{chap:lit} reviews the literature on irrigation control,
optimal and predictive control, reinforcement learning, spatial flow
modeling, and the supporting sensing and communication infrastructure;
the chapter concludes by identifying the specific research gap
addressed by the thesis. Chapter~\ref{chap:system} presents the
agent-based virtual plant: the study site and digital elevation model,
the directed flow graph governing lateral water exchange, the climate
data and scenario selection, the soil and crop parameterization, the
state-space dynamics, the action space and water-budget constraint,
and the unified software framework that ensures fair comparison
across controllers. Chapter~\ref{chap:controllers} introduces the four
controller architectures and develops the constrained MPC formulation
in detail, including the multi-objective cost function with economic
anchoring, the symbolic control-oriented model, the systematic
weight-sensitivity analysis, and the SAC architecture with its CTDE
training procedure. Chapter~\ref{chap:results} presents the
comparative results across the 9-cell factorial grid of climate
scenario and budget allocation. Chapter~\ref{chap:implementation}
addresses the considerations relevant to physical deployment.
Chapter~\ref{chap:conclusion} summarizes the principal findings,
identifies the limitations of the present work, and proposes
directions for future research.
